\documentclass[11pt]{article}
\usepackage{amsmath,amssymb,amsthm,hyperref}
\newtheorem{theorem}{Theorem}[section]
\newtheorem{lemma}[theorem]{Lemma}
\newtheorem{remark}[theorem]{Remark}
\newcommand{\CH}{\operatorname{CH}}
\newcommand{\ch}{\operatorname{ch}}
\title{Integral projection correction and the singular-branch boundary of spectral sources}
\date{28 October 2026; base checkpoint 02bba4d}
\begin{document}\maketitle
\section{Contract and three executed routes}
The original input is an arbitrary smooth projective complex $X$, codimension
$p$, and prescribed $\alpha\in H^{2p}(X,\mathbb Q)\cap H^{p,p}(X)$.
The required output remains an actual rational Chow class on the same $X$
with canonical Betti class $\alpha$. None of the following identities assumes
that this output exists. Let $q:J\to X$ be the actual affine Jouanolou bundle.
The alpha-dependent polynomial near-projectors constructed in the October 24
record are retained, including their local topological class on a compact
cohomology-detecting core. The new tests are: (i) finite nilpotent exactification;
(ii) a repeated anti-eigenline square across smooth ramification, beyond the
previous unramified calculation; (iii) independent anti-eigenlines on a
connected finite etale cover of the whole original variety.

\section{Integral finite correction in every associative ring}
For an element $a$ of any unital associative ring put
\[
 e(a)=a^2-a,\qquad T(a)=a-(2a-1)e(a)=3a^2-2a^3.
\]
\begin{theorem}[Exact defect doubling, with no inverse of two]
One has $e(T(a))=e(a)^2(4e(a)-3)$. If $e(a)^{2^k}=0$, then
$T^k(a)$ is an idempotent. In particular every nilpotent defect admits a finite
integer-coefficient polynomial correction.
\end{theorem}
\begin{proof}
Expand the displayed polynomials. All factors are polynomials in the same
element, so they commute even in a noncommutative ambient ring. Consequently
$e(T(a))^m=e(a)^{2m}(4e(a)-3)^m$. Induct on $k$.
If $e(a)^n=0$, choose $k$ with $2^k\ge n$; a larger vanishing power follows
by multiplication. Thus $\lceil\log_2 n\rceil$ rounds suffice for $n\ge1$.
The existence statement includes the zero ring; no division is used.
\end{proof}
This is the familiar idempotent-lifting mechanism, not a global novelty claim.
The branch previously recorded only first-order projector conjugacy; here the
algorithm and stopping criterion are integral and at arbitrary finite depth.

\begin{theorem}[Explicit conjugacy from a nilpotent difference]
Let $p^2=p$, $s^2=s$, and $(s-p)^n=0$ in any associative ring. Set
\[
 u=sp+(1-s)(1-p),\quad v=ps+(1-p)(1-s),\quad z=1-(s-p)^2.
\]
Then $up=su$, $uv=vu=z$, $u$ is a unit, and $s=upu^{-1}$.
\end{theorem}
\begin{proof}
Expand using only the two idempotence relations. Since $(s-p)^2$ is
nilpotent, $z^{-1}=\sum_{j=0}^{n-1}(s-p)^{2j}$ is a two-sided inverse.
The elements $vz^{-1}$ and $z^{-1}v$ are respectively a right and left inverse
of $u$. They coincide: if $Lu=1=uR$, then $R=(Lu)R=L(uR)=L$.
Multiply $up=su$ by $u^{-1}$. Thus image projective modules are isomorphic.
In particular a nilpotent thickening around a fixed algebraic source cannot
change its $K_0$ class. This does not identify a general small analytic error
with a nilpotent algebraic error.
\end{proof}

\section{A torsion-line calculation without an immersion GRR assumption}
\begin{lemma}[Nilpotent involutions]
If $x^2=1$ in a commutative ring and $b=x-1$, then
$b^{j+1}=(-2)^j b$ for every $j\ge0$. Thus $b^{j+1}=0$ implies
$2^j b=0$; over a rational algebra $x=1$.
\end{lemma}
\begin{proof}
$b^2=-2b$ and induction give the formula. The rational conclusion uses the
inverse of the nonzero scalar $(-2)^j$, not cancellation in an arbitrary ring.
\end{proof}
For a line bundle $K$ on a smooth finite-dimensional quasi-projective scheme,
$[K]-1$ is nilpotent in algebraic $K_0$. Here is a sufficient direct argument.
Choose an ample twist $A$ so that $A$ and $K\otimes A$ are generated by finitely
many global sections. The exact Koszul complexes of generating sections give
$(1-[A]^{-1})^N=0$ and $(1-[K\otimes A]^{-1})^M=0$. These are commuting
nilpotent elements, and their difference and multiplication by a line-class
unit show $[K]-1$ nilpotent. Thus $K^{\otimes2}\simeq O$ implies
$[K]=1$ in $K_0\otimes\mathbb Q$.

\section{The repeated anti-source is concentrated over singular branching}
Let $J$ be smooth and integral over $\mathbb C$, let $b$ be a regular traceless $2\times2$
matrix, and write $b^2=dI$. Assume $d$ is not identically zero and its zero
divisor $D$ is reduced. Let $S=\operatorname{Sing}(D)$ and $U=J\setminus S$.
On $U$, let $\pi:Y=\{t^2=d\}\to U$ be the finite double cover with involution
$\sigma(t)=-t$. It is smooth; the ramification divisor $R=\{t=0\}$ is smooth
and principal. The eigenline bundles are the actual kernels
$L_+=\ker(b-t)$, $L_-=\ker(b+t)$ in $O_Y^2$.

\begin{theorem}[Smooth-branch neutralization]
In $K_0(Y)\otimes\mathbb Q$ one has
\[
 [L_+]+[L_-]=2,\qquad [L_+][L_-]=1,\qquad
 ([L_+]-[L_-])^r=0\quad(r\ge2).
\]
\end{theorem}
\begin{proof}
At a point of $D$, the matrix $b$ is nonzero: if all its entries vanished,
the differential of $d=b_{11}^2+b_{12}b_{21}$ would vanish too, contrary to
smoothness. It is rank-one nilpotent on $D$. Thus
$K=\ker(b|_D)=\operatorname{im}(b|_D)$ is a line bundle and
$0\to K\to O_D^2\to K\to0$ gives $K^{\otimes2}\simeq O_D$.
Locally choose $w$ with $(bw,w)$ a basis; then $b$ has matrix
$\left(\begin{smallmatrix}0&1\\d&0\end{smallmatrix}\right)$ and the two
eigenlines are generated by $(1,t)$ and $(1,-t)$. The addition map gives
the exact sequence
\[
 0\to L_+\oplus L_-\to O_Y^2\to i_*K\to0,
 \qquad i:R\simeq D\hookrightarrow Y.
\]
Its determinant vanishes simply on $R$, so
$L_+\otimes L_-=O_Y(-R)\simeq O_Y$.
The torsion-line lemma says $[K]=1$ rationally; exact pushforward from $R$
then gives $[i_*K]=[i_*O_R]=0$, since the principal-divisor sequence is
$0\to O_Y\xrightarrow{t}O_Y\to O_R\to0$.
The sum relation follows. In any commutative ring with $x+y=2$ and $xy=1$,
$(x-y)^2=(x+y)^2-4xy=0$, and all higher powers vanish.
No Grothendieck--Riemann--Roch formula for a closed immersion is needed here.
\end{proof}

\begin{theorem}[Codimension-two source survival requires codimension-two singular support]
Choose a proper equivariant resolution $\widetilde Y\to\{t^2=d\}$ which is
an isomorphism over $U$, and extensions of the two eigenlines which agree
with their actual kernels over $U$. For $r\ge2$ let
$\Delta=[\bar L_+]-[\bar L_-]$ and
$\xi_r=\pi_*(\Delta^r)\in K_0(J)\otimes\mathbb Q$.
Every $\ch_p(\xi_r)$ restricts to zero on $U$. It is in the image of
$\CH_{\dim J-p}(S)_\mathbb Q\to\CH^p(J)_\mathbb Q$.
If $\operatorname{codim}_J S>p$, this class is zero.
In particular the quadratic source in codimension two vanishes whenever
$\operatorname{codim}_J S\ge3$.
\end{theorem}
\begin{proof}
On the unchanged smooth open, $\Delta^r=0$ rationally by the theorem.
Restriction commutes with proper coherent pushforward and Chern character.
Apply the Chow localization exact sequence for $S\subset J$.
When $\dim S<\dim J-p$, there are no cycles of the required dimension on $S$.
For the actual Jouanolou $J\to X$, its true Chow homotopy inverse returns
zero on the original $X$. This says nothing about arbitrary new supports.
\end{proof}
This sharpens the older boundary localization from the entire discriminant
to its singular locus, under the displayed reducedness and kernel-extension
conditions. If $d$ has repeated factors or arbitrary boundary twists are
inserted, these hypotheses must be rechecked; those cases are not suppressed.
For a map $J\to\mathbb A^3$, $b=\left(\begin{smallmatrix}x&y\\z&-x\end{smallmatrix}\right)$,
which is transverse to the smooth cone $x^2+yz=0$ and whose inverse image of
the vertex has codimension three, the singular-locus condition holds. Existence
of an alpha-matched nonzero class from deliberately nontransverse maps is not
implied. This is the surviving geometric search, not a new assumed support law.

\section{A constructive singular-stratum extraction by three twists}
The preceding obstruction specifies where a nonzero repeated source must live.
Here is an actual geometric return from that surviving locus. Suppose a smooth
codimension-two stratum $S\subset J$ has a transverse $A_n$ surface-resolution
chain in the proper spectral resolution, with globally defined exceptional
divisors $E_1,\ldots,E_n$. Assume the deck exchanges $E_i$ and $E_j$, $i\ne j$,
and the transverse intersection equations are
\[
 \pi_*(E_i^2)=\pi_*(E_j^2)=-2[S],\qquad
 \pi_*(E_iE_j)=\epsilon[S],\quad
 \epsilon=\begin{cases}1&|i-j|=1,\\0&|i-j|>1.\end{cases}
\]
These equations are required on all of the chosen stratum; additional singular
strata or monodromy of the exceptional components require separate terms.
Put $D=E_i-E_j$ and $c=-4-2\epsilon\in\{-4,-6\}$. Thus
$\sigma^*D=-D$ and $\pi_*(D^2)=c[S]$.
For any equivariantly exchanged line extensions $\bar L_\pm$, put
$a=c_1(\bar L_+)-c_1(\bar L_-)$ and form the actual Chow cycles
\[
 Z(k)=(q^*)^{-1}\pi_*\bigl(a+2kD\bigr)^2\in\CH^2(X)_\mathbb Q,
 \qquad k\in\mathbb Z.
\]
These are genuine intersections of algebraic divisor classes followed by
proper Chow pushforward, not merely formal period variables. Twisting
$\bar L_\pm$ by $O(\pm kD)$ gives exactly the displayed first-Chern difference.
\begin{theorem}[The starting eigenline cancels]
The three constructed cycles satisfy the exact original-$X$ Chow equality
\[
 \boxed{\frac{Z(1)+Z(-1)-2Z(0)}{8c}=(q^*)^{-1}[S].}
\]
For the two exchanged components of an $A_2$ chain, $c=-6$, so the coefficient
is $-1/48$. For a nonadjacent exchanged pair it is $-1/32$.
\end{theorem}
\begin{proof}
Expand $(a+2D)^2+(a-2D)^2-2a^2=8D^2$ in the actual Chow ring.
Push forward and use the stated transverse intersection equations; apply
the affine-bundle Chow inverse. The unknown horizontal class $a$ and its mixed
term cancel exactly. Canonical Betti equality follows functorially.
\end{proof}
This is an explicit extraction from independently present singular geometry;
it does not assert such a stratum matches every prescribed $\alpha$. An $A_1$
chain has no exchanged distinct pair. Furthermore the tempting universal
choice $b=\left(\begin{smallmatrix}f&g\\g^2&-f\end{smallmatrix}\right)$,
with global regular functions $f,g$, gives a transverse $A_2$ locus
$S=V(f,g)$ when $df,dg$ are independent. But this global regular complete
intersection has $[S]=c_2(O_J^2)=0$ in $\CH^2(J)$, so its extracted original
source is zero. A target-sensitive surviving construction must change this
support geometry (or separate nontrivial components), not just add twists.

\section{Connected whole-cover rigidity for independent mixed sources}
\begin{theorem}[Projective positivity detects the connectedness restriction]
Let $f:X'\to X$ be a connected finite etale double cover of connected smooth
projective complex varieties, with deck $\sigma$, and let $J'=J\times_X X'$.
Suppose an algebraic line bundle $L$ on $J'$ satisfies
$L\oplus\sigma^*L\simeq O_{J'}^2$. Then its canonical rational Betti first
Chern class is zero. Every mixed product of such first-Chern differences
therefore has zero canonical Betti class.
\end{theorem}
\begin{proof}
Affine-bundle Chow invariance lifts $c_1(L)$ to a rational divisor class on
$X'$ and is injective on its square. Write its Betti class as $l$.
Complementarity and determinant give $\sigma^*l=-l$ and $l^2=0$.
Let $h=f^*h_X$ be an ample class. Since $X'$ is connected, integration is
one scalar and $\int l h^{n-1}=\int\sigma^*l h^{n-1}=-\int l h^{n-1}=0$.
For $n\ge2$ the divisor Hodge index theorem is negative definite on this
primitive rational $(1,1)$ subspace. Its square pairing is zero, so $l=0$.
Concretely, Theorem 18.15 of Christian Schnell's Fall 2024 Class 18 notes
gives $-\int l^2h^{n-2}>0$ for every nonzero real primitive $(1,1)$ class;
this is the $k=2$, $p=q=1$ specialization of Hodge--Riemann.
For $n=1$, degree determines $H^2$, giving the same conclusion.
The assertion is Betti vanishing; it does not assert $L$ trivial or
$c_1^{\CH}(L)=0$ in Chow. On a split double cover integration has two
component values exchanged by $\sigma$, so the argument fails. The existing
nonzero $(\mathbb P^1)^r$ split-cover mixed-source example is preserved.
\end{proof}
The theorem is stated for covers already descended to $X$. Riemann existence
and the contractible fibers of $J^{an}\to X^{an}$ also descend finite etale
covers, but this comparison is not required as an unproved Lean instance.

\section{Evidence and remaining mathematical distinction}
The Lean drafts encode integral defect doubling, the explicit unit and
conjugacy, the torsion-line algebra, exact boundary-defect identities and
the primitive-positivity implication. They do not encode schemes, resolution,
Chow localization or the geometric Hodge index theorem. Hand proofs above
separate those inputs. Computations test free-ring identities, finite exact
nilpotent quotients including characteristics two and three, and a ramification
normal form. None constructs a canonical Chow representative for every alpha.
The first-order tangent and full spectral-trace mechanisms were already in
the branch; they are not counted as new geometry here.
Primary comparisons checked on 28 October 2026:
\href{https://stacks.math.columbia.edu/tag/00J9}{Stacks 00J9}
(commutative idempotent lifting; the noncommutative extension above is proved
directly), \href{https://stacks.math.columbia.edu/tag/02RX}{Stacks 02RX}
(Chow localization), and
\href{https://www.math.stonybrook.edu/~cschnell/mat545/class18.pdf}{Schnell,
Theorem 18.15, p.~6} (the exact positivity statement).
\end{document}
